\begin{description}
\item[(T06.1)]
  % FIGURE NEEDED: schematic top view of the Galaxy (2025_theory_solution.pdf
  % p.17): Sun on its orbit with longitude directions l = 0, 45, 90, 135, 180,
  % 270 deg marked as arrows, GC at centre, dotted circular orbits, and the
  % circle Omega = Omega_0 through the Sun.
  The concept of differential rotation of the Galaxy is used here.
  \begin{itemize}
  \item For $l = 45^\circ$:
    \begin{itemize}
    \item From Sun to the GC, the quantity $(\Omega - \Omega_0)$ increases,
      which implies $V_{\mathrm{r}}$ increases.
    \item As shown in the figure, the line of sight intersects circular orbits
      that are at varying distances from the GC.
    \item The line of sight is closest to the GC when
      $D = R_0\cos l = 6.01$\,kpc. $V_{\mathrm{r}}$ reaches the maximum value
      here.
    \item Beyond this, $(\Omega - \Omega_0)$ decreases implying
      $V_{\mathrm{r}}$ decreases.
    \item $V_{\mathrm{r}} = 0$ when the line of sight intersects the Sun's
      orbit. Here, $D = 2R_0\cos l = 12.02$\,kpc
    \item Beyond the Sun's orbit, $\Omega < \Omega_0$. $V_{\mathrm{r}}$ is
      negative and its magnitude increases with increasing $D$.
    \end{itemize}
  \item $l = 135^\circ$
    \begin{itemize}
    \item The line of sight lies outside the Sun's orbit. Here,
      $V_{\mathrm{r}}$ is negative, there is no peak, and its magnitude
      increases with increasing $D$.
    \end{itemize}
  \end{itemize}
  The plot of the line-of-sight velocity as a function of the distance from
  the Sun for two values of $l$ is as follows:
  % FIGURE NEEDED: sketch of V_r vs D from 0 to 2R_0 (2025_theory_solution.pdf
  % p.18): l = 45 deg curve starts at (0,0), rises to a single positive
  % maximum at D = 6 kpc, crosses zero at D = 12 kpc and is negative beyond;
  % l = 135 deg curve starts at (0,0) and is negative throughout, decreasing
  % monotonically.
\item[(T06.2)] Consider $l = 90$,
  \begin{align*}
    V_{\mathrm{r}} &= 0 \text{ implying } \Omega = \Omega_0
       \tag*{(0.5)}\\
    V_{\mathrm{t}} &= -\Omega_0 D = -2.75\,\mathrm{km\,s^{-1}}
       \tag*{(0.5)}\\
    \Omega_0 &= 8.91\times10^{-16}\,\mathrm{rad\,s^{-1}}
       \tag*{(1.0)}\\
    P &= \frac{2\pi}{\Omega_0} = 224\,\mathrm{Myr}
       \tag*{(1.0)}
  \end{align*}
\item[(T06.3a)] \mbox{}
  \begin{description}
  \item[(I)] The observed rotation curve is flat near the Sun.

    $V$ is constant $\rightarrow$
    $\left.\dfrac{dV}{dR}\right|_{R=R_0} = 0$. \hfill(1.0)
  \item[(II)] If the whole mass of the Galaxy is assumed to be concentrated
    at the centre, then we consider Keplerian motion.
    \[
      V = \sqrt{\frac{GM}{R}}
      \;\rightarrow\;
      \left.\frac{dV}{dR}\right|_{R=R_0} = -\frac{1}{2}\frac{V_0}{R_0}.
    \]
    \hfill(1.0)
  \end{description}
\item[(T06.3b)] In the solar neighbourhood, $(\Omega - \Omega_0)$ will be
  very small. Using the approximation
  \[
    f(x) \simeq f(x_0)
      + \left.\frac{df}{dx}\right|_{x=x_0}(x - x_0),
    \text{ for } x \approx x_0, \text{ given in the Data Sheet,}
  \]
  \begin{align*}
    (\Omega - \Omega_0)
      &= \left.\frac{d\Omega}{dR}\right|_{R=R_0}(R - R_0)
       \tag*{(2.0)}\\
      &= \left[\left(\frac{1}{R}\frac{dV}{dR}\right)\bigg|_{R=R_0}
         - \left(\frac{1}{R^2}V\right)\bigg|_{R=R_0}\right](R - R_0)\\
      &= \frac{1}{R_0^2}\left[R_0\left.\frac{dV}{dR}\right|_{R=R_0}
         - V_0\right](R - R_0)
  \end{align*}
  Thus the line-of-sight velocity expression becomes
  \[
    V_{\mathrm{r}} = \frac{1}{R_0^2}
      \left[R_0\left.\frac{dV}{dR}\right|_{R=R_0} - V_0\right]
      (R - R_0)R_0\sin l
  \]
  Further, in the solar neighbourhood, $D \ll R_0$. So we can approximate
  \[
    (R - R_0) \approx -D\cos l
  \]
  \hfill(1.0)

  Thus, the line-of-sight velocity can be expressed as,
  \begin{align*}
    V_{\mathrm{r}}
      &= \frac{1}{R_0^2}
         \left[R_0\left.\frac{dV}{dR}\right|_{R=R_0} - V_0\right]
         (-D\cos l)R_0\sin l\\
      &= \frac{1}{2}
         \left[\frac{V_0}{R_0} - \left.\frac{dV}{dR}\right|_{R=R_0}\right]
         D\sin 2l
  \end{align*}
  Comparing with $V_{\mathrm{r}} = AD\sin 2l$, we get
  \[
    A = \frac{1}{2}
        \left[\frac{V_0}{R_0} - \left.\frac{dV}{dR}\right|_{R=R_0}\right]
  \]
  \hfill(2.0)

  Also, in the solar neighbourhood, $\Omega D \approx \Omega_0 D$. Using this
  approximation and $(R - R_0) \approx -D\cos l$, the tangential velocity is
  \hfill(1.0)
  \begin{align*}
    V_{\mathrm{t}} &= (\Omega - \Omega_0)R_0\cos l - \Omega D\\
      &= \frac{1}{R_0^2}
         \left[R_0\left.\frac{dV}{dR}\right|_{R=R_0} - V_0\right]
         (-D\cos l)R_0\cos l - \Omega_0 D\\
      &= \left[\frac{V_0}{R_0}
         - \left.\frac{dV}{dR}\right|_{R=R_0}\right]D\cos^2 l - \Omega_0 D\\
      &= \frac{1}{2}\left[\frac{V_0}{R_0}
         - \left.\frac{dV}{dR}\right|_{R=R_0}\right]D(1 + \cos 2l)
         - \frac{V_0}{R_0}D\\
      &= AD\cos 2l
         + \frac{1}{2}\left[\frac{V_0}{R_0}
         - \left.\frac{dV}{dR}\right|_{R=R_0}\right]D - \frac{V_0}{R_0}D\\
      &= AD\cos 2l
         - \frac{1}{2}\left[\frac{V_0}{R_0}
         + \left.\frac{dV}{dR}\right|_{R=R_0}\right]D
  \end{align*}
  Comparing with $V_{\mathrm{t}} = AD\cos 2l + BD$, we get
  \[
    B = -\frac{1}{2}
        \left[\frac{V_0}{R_0} + \left.\frac{dV}{dR}\right|_{R=R_0}\right]
  \]
  \hfill(2.0)
\item[(T06.3c)] For a flat rotation curve
  $\left(\dfrac{dV}{dR} = 0\right)$, $F_{\mathrm{I}} = -1$. \hfill(1.0)

  For Keplerian rotation
  $\left(\dfrac{dV}{dR} = -\dfrac{1}{2}\dfrac{V}{R}\right)$,
  $F_{\mathrm{II}} = -3$. \hfill(1.0)
\end{description}
